%% kapa: exam=1 round=30 year=2019 period=2 form=A subject=accounting subject_ko=회계학 seq=36 number=76
%% source: Q-Net 감정평가사 기출문제 (cst00309) artlSeq=5205205 file=extracted/1008/5205205/2019년 제30회 감정평가사 1차시험 문제지/2019년 제30회 감정평가사 1차시험 문제지/2019년도 제30회 감정평가사 1차 2교시 A형.hwp
%% conversion: kordoc 4.24.51 -> tools/gen1.py (auto); review=proofread-sonnet; flags=none
\begin{question}{76}
표준원가계산에 관한 설명으로 옳은 것을 모두 고른 것은?
\begin{framedtext}
ㄱ. 표준원가계산제도는 전부원가계산에서 적용할 수 있으나 변동원가계산에서는 적용할 수 없다. \par
ㄴ. 표준원가계산제도는 종합원가계산제도에 적용이 가능하다. \par
ㄷ. 직접재료원가 가격차이를 구입시점에서 분리하든 사용시점에서 분리하든 직접재료원가 능률차이는 동일하다. \par
ㄹ. 고정제조간접원가의 예산차이는 실제투입량 변동예산과 실제산출량 변동예산의 차이를 의미한다. \par
\end{framedtext}
\begin{choices}
  \item ㄱ, ㄴ
  \item ㄱ, ㄷ
  \item ㄴ, ㄷ
  \item ㄴ, ㄹ
  \item ㄷ, ㄹ
\end{choices}
\end{question}
\answer{3}
